%% GENERATED by python/paper/build.py from results/ -- do not edit by hand.
%% Rebuild:  .venv\Scripts\python.exe python\paper\build.py --only OECD_InequalityMeasures
%% Source:   data/oecdFigure1.csv, oecd29_distribution.csv, oecd26_theta.csv
\textit{Note:} The 21 countries of figure \ref{fig:US:OECDdata}. The Gini is that of disposable income (WIID); high over low income, $y_3/y_1$, and the dispersion $\tilde V = \sum_i \gamma_i (y_i - 1)^2/y_i$ are of the three income groups the model is fed, at the U.S.\ shares of the population (58\%, 18\%, 24\%, from the bottom), in the gross earnings of men working full time (OECD). The solid line is the least-squares fit over all countries, the dashed one over all but the U.S.; each panel gives Pearson\textquotesingle s $r$ and Spearman\textquotesingle s rank correlation, with the U.S.\ and, in parentheses, without it (table \ref{table:OA:OECD_InequalityCorrelations}).
